\documentclass{article}

% packages
\usepackage{amsmath, amsthm, thmtools, amsfonts, amssymb, luacode, catchfile, hyperref, ifthen, graphicx, bbm}
\ifcsname c@kobocompile\endcsname
	\usepackage[a5paper, total={1072pt, 1448pt}, margin=10pt, includeheadfoot]{geometry} % set page margins
\else
	\usepackage[a4paper, margin=50pt, includeheadfoot]{geometry}
\fi
\usepackage[shortlabels]{enumitem}
\usepackage[skip=3pt, indent=0pt]{parskip}

% language
\usepackage[bidi=basic, layout=tabular, provide=*]{babel}
\ifcsname c@english\endcsname
	\babelprovide[main, import]{english}
\else
	\babelprovide[main, import, alph=letters.gershayim]{hebrew}
	\setlist[enumerate,2]{label={\alph*.}}
	\babelprovide{rl}
\fi
\babelfont{rm}[Renderer=Harfbuzz]{Libertinus Serif}
\babelfont{sf}{Libertinus Sans}
\babelfont{tt}{Libertinus Mono}

% style
\ifcsname c@stoppagebreaks\endcsname
\else
	\AddToHook{cmd/section/before}{\clearpage}	% Add line break before section
\fi
\linespread{1.3}
\setcounter{secnumdepth}{0}		% Remove default number tags from sections, this won't do well with theorems
\AtBeginDocument{\setlength{\belowdisplayskip}{3pt}}
\AtBeginDocument{\setlength{\abovedisplayskip}{3pt}}
\graphicspath{ {../images/} }

% operators
\DeclareMathOperator\cis{cis}
\DeclareMathOperator\Sp{Span}
\DeclareMathOperator\tr{tr}
\DeclareMathOperator\im{Im}
\DeclareMathOperator\re{Re}
\DeclareMathOperator\diag{diag}
\DeclareMathOperator*\lowlim{\underline{lim}}
\DeclareMathOperator*\uplim{\overline{lim}}
\DeclareMathOperator\rng{rng}
\DeclareMathOperator\Sym{Sym}
\DeclareMathOperator\Arg{Arg}
\DeclareMathOperator\Log{Log}
\DeclareMathOperator\dom{dom}
\DeclareMathOperator\supp{Supp}
\DeclareMathOperator\var{Var}
\DeclareMathOperator\cov{Cov}
\DeclareMathOperator\End{End}
\DeclareMathOperator\id{id}
\DeclareMathOperator\cl{cl}
\DeclareMathOperator\gal{Gal}
\DeclareMathOperator\dist{dist}
\DeclareMathOperator\form{form}
\DeclareMathOperator\term{term}
\DeclareMathOperator\sent{sent}
\DeclareMathOperator\Th{Th}
\DeclareMathOperator\Hom{Hom}
\DeclareMathOperator\Aut{Aut}
\DeclareMathOperator\Char{char}
\DeclareMathOperator\vol{vol}
\DeclareMathOperator\GL{GL}
\DeclareMathOperator\SL{SL}
\DeclareMathOperator\lcm{lcm}

% commands
\newcommand{\cupdot}{\mathbin{\mathaccent\cdot\cup}}
\newcommand{\indicator}{\mathbbm{1}}
\newcommand{\AAA}[0]{\mathbb{A}}
\newcommand{\Aa}[0]{\mathcal{A}}
\newcommand{\BB}[0]{\mathbb{B}}
\newcommand{\Bb}[0]{\mathcal{B}}
\newcommand{\CC}[0]{\mathbb{C}}
\newcommand{\Cc}[0]{\mathcal{C}}
\newcommand{\DD}[0]{\mathbb{D}}
\newcommand{\Dd}[0]{\mathcal{D}}
\newcommand{\EE}[0]{\mathbb{E}}
\newcommand{\FF}[0]{\mathbb{F}}
\newcommand{\Ff}[0]{\mathcal{F}}
\newcommand{\Ii}[0]{\mathcal{I}}
\newcommand{\Gg}[0]{\mathcal{G}}
\newcommand{\HH}[0]{\mathbb{H}}
\newcommand{\Hh}[0]{\mathcal{H}}
\newcommand{\Kk}[0]{\mathcal{K}}
\newcommand{\Ll}[0]{\mathcal{L}}
\newcommand{\Mm}[0]{\mathcal{M}}
\newcommand{\NN}[0]{\mathbb{N}}
\newcommand{\Nn}[0]{\mathcal{N}}
\newcommand{\PP}[0]{\mathbb{P}}
\newcommand{\Pp}[0]{\mathcal{P}}
\newcommand{\QQ}[0]{\mathbb{Q}}
\newcommand{\RR}[0]{\mathbb{R}}
\newcommand{\Rr}[0]{\mathcal{R}}
\newcommand{\Ss}[0]{\mathcal{S}}
\newcommand{\TT}[0]{\mathbb{T}}
\newcommand{\Tt}[0]{\mathcal{T}}
\newcommand{\UU}[0]{\mathbb{U}}
\newcommand{\Uu}[0]{\mathcal{U}}
\newcommand{\Vv}[0]{\mathcal{V}}
\newcommand{\Ww}[0]{\mathcal{W}}
\newcommand{\Xx}[0]{\mathcal{X}}
\newcommand{\Yy}[0]{\mathcal{Y}}
\newcommand{\ZZ}[0]{\mathbb{Z}}
\newcommand{\Zz}[0]{\mathcal{Z}}
\newcommand{\acts}[0]{\circlearrowright}
\newcommand{\explain}[2] {
	\begin{flalign*}
		 && \text{#2} && \text{#1}
	\end{flalign*}
}
\newcommand{\maketitleprint}[1][green]{
	\begin{center}
		\includegraphics[width=6cm]{dragon_#1}
	\end{center}
}

% theorem commands
\newtheoremstyle{c_remark}
	{}	% Space above
	{}	% Space below
	{}% Body font
	{}	% Indent amount
	{\bfseries}	% Theorem head font
	{}	% Punctuation after theorem head
	{.5em}	% Space after theorem head
	{\thmname{#1}\thmnumber{ #2}\thmnote{ \normalfont{\text{(#3)}}}}	% head content
\newtheoremstyle{c_definition}
	{3pt}	% Space above
	{3pt}	% Space below
	{}% Body font
	{}	% Indent amount
	{\bfseries}	% Theorem head font
	{}	% Punctuation after theorem head
	{.5em}	% Space after theorem head
	{\thmname{#1}\thmnumber{ #2}\thmnote{ \normalfont{\text{(#3)}}}}	% head content
\newtheoremstyle{c_plain}
	{3pt}	% Space above
	{3pt}	% Space below
	{\itshape}% Body font
	{}	% Indent amount
	{\bfseries}	% Theorem head font
	{}	% Punctuation after theorem head
	{.5em}	% Space after theorem head
	{\thmname{#1}\thmnumber{ #2}\thmnote{ \text{(#3)}}}	% head content

\ifcsname c@english\endcsname
	\theoremstyle{plain}
	\newtheorem{theorem}{Theorem}[section]
	\newtheorem*{theorem*}{Theorem}
	\newtheorem{lemma}[theorem]{Lemma}
	\newtheorem{proposition}[theorem]{Proposition}
	\newtheorem*{proposition*}{Proposition}

	\theoremstyle{definition}
	\newtheorem{definition}[theorem]{Definition}
	\newtheorem*{definition*}{Definition}
	\newtheorem{example}{Example}[section]
	\newtheorem{exercise}{Exercise}[section]

	\theoremstyle{remark}
	\newtheorem*{remark}{Remark}
	\newtheorem*{solution}{Solution}
	\newtheorem{corollary}[theorem]{Corollary}
	\newtheorem{notation}[theorem]{Notation}
\else
	\theoremstyle{c_plain}
	\newtheorem{theorem}{משפט}[section]
	\newtheorem*{theorem*}{משפט}
	\newtheorem{lemma}[theorem]{למה}
	\newtheorem{proposition}[theorem]{טענה}
	\newtheorem*{proposition*}{טענה}

	\theoremstyle{c_definition}
	\newtheorem{definition}[theorem]{הגדרה}
	\newtheorem*{definition*}{הגדרה}
	\newtheorem{example}{דוגמה}[section]
	\newtheorem{exercise}{תרגיל}[section]

	\theoremstyle{c_remark}
	\newtheorem*{remark}{הערה}
	\newtheorem*{solution}{פתרון}
	\newtheorem{corollary}[theorem]{מסקנה}
	\newtheorem{notation}[theorem]{סימון}
\fi

% Questions related commands
\newcounter{question}
\setcounter{question}{1}
\newcounter{sub_question}
\setcounter{sub_question}{1}

\ifcsname c@english\endcsname
	\newcommand{\question}[1][0]{
		\ifthenelse{#1 = 0}{}{\setcounter{question}{#1}}
		\section{Question \arabic{question}}
		\addtocounter{question}{1}
		\setcounter{sub_question}{1}
	}

	\newcommand{\subquestion}[1][0]{
		\ifthenelse{#1 = 0}{}{\setcounter{sub_question}{#1}}
		\subsection{Part \alph{sub_question}}
		\addtocounter{sub_question}{1}
	}
\else
	\newcommand{\question}[1][0]{
		\ifthenelse{#1 = 0}{}{\setcounter{question}{#1}}
		\section{שאלה \arabic{question}}
		\addtocounter{question}{1}
		\setcounter{sub_question}{1}
	}

	\newcommand{\subquestion}[1][0]{
		\ifthenelse{#1 = 0}{}{\setcounter{sub_question}{#1}}
		\subsection{סעיף \localecounter{letters.gershayim}{sub_question}}
		\addtocounter{sub_question}{1}
	}
\fi

% import lua and start of document
\directlua{common = require ('../common')}

\GetEnv{AUTHOR}

% headers
\author{\AUTHOR}
\date\today

\title{פתרון מטלה 05 --- יסודות באלגברה קומוטטיבית}

\DeclareMathOperator{\Ann}{Ann}

\begin{document}
\maketitle
\maketitleprint[red]

\question[2]{}
יהיו $M, M', M'', N, N', N''$ מודולים מעל $R$.
נניח כי שתי הסדרות,
\[
	M'
	\xrightarrow{f} M
	\xrightarrow{g} M''
	\to 0,
	\qquad
	N'
	\xrightarrow{\varphi} N
	\xrightarrow{\psi} N''
	\to 0
\]
הן מדויקות.
נגדיר את,
\[
	\theta : (M \otimes_R N') \oplus (M' \otimes_R N) \to M \otimes_R N,
	\qquad
	\theta(x, y)
	= \id \otimes \varphi(x) + f \otimes \id(y)
\]
ונראה כי הסדרה,
\[
	(M \otimes_R N') \oplus (M' \otimes_R N)
	\xrightarrow{\theta} M \otimes_R N
	\xrightarrow{g \otimes \psi} M'' \otimes_R N''
	\to 0
\]
היא סדרה מדויקת.
\begin{proof}
	שקול להוכחת דיוק מימין.
\end{proof}

\question{}
יהי $f : A \to B$ הומומורפיזם חוגים, ויהי $N$ מודול מעל $B$, על־ידי צמצום סקלרים $N$ הוא גם מודול מעל $A$.
נסמן $N_B = B \otimes_A N$, זהו מודול מעל $B$, ועל־ידי הרחבת סקלרים נגדיר הומומורפיזם,
\[
	\varphi : N \to N_B,
	\qquad
	\varphi(m) = 1 \otimes m
\]
נראה ש־$\varphi$ חד־חד ערכית.
\begin{proof}
	נגדיר את ההעתקה $\psi : B \times N \to N$ על־ידי $\psi(b, n) = b \cdot n$.
	נראה כי זוהי העקה בילינארית מעל $A$,
	\[
		a . \psi(b, n)
		= f(a) \cdot b n
		= (f(a) \cdot b) \cdot n
		= \psi(f(a) b, n)
	\]
	אבל גם,
	\[
		a . \psi(b, n)
		= b f(a) n
		= \psi(b, f(a) n)
	\]
	ולכן זוהי העתקה $A$־בילינארית, ונסיק כי קיימת ההרמה,
	\[
		\tilde{\psi} : B \otimes_A N \to N,
		\qquad
		\tilde{\psi}(b \otimes n)
		= b n
	\]
	כלומר $\tilde{\psi} : N_B \to N$ הומומורפיזם.

	יהי $x \in N$ כך ש־$\varphi(x) = 0$, כלומר $1 \otimes x = 0$, אז,
	\[
		0
		= \tilde{\psi}(0)
		= \tilde{\psi}(1 \otimes x)
		= 1 \cdot x
	\]
	כלומר $x = 0$ כפי שרצינו.
\end{proof}

\question{}
יהי $f : A \to B$ הומומורפיזם חוגים, ויהי $N$ מודול מעל $B$, על־ידי צמצום סקלרים גם מודול מעל $A$.
נסמן $N_B = B \otimes_A N$, וזהו $B$־מודול על־ידי הרחבת סקלרים.
נגדיר את,
\[
	\varphi : N \to N_B,
	\qquad
	n \mapsto 1 \otimes n
\]
ונבחין כי זהו הומומורפיזם.
נראה כי,
\[
	0
	\to N
	\xrightarrow{\varphi} N_B
	\xrightarrow{\pi_{\im \varphi}} N_B / \im \varphi
	\to 0
\]
היא סדרה מדויקת ומתפצלת.
\begin{proof}
	$\varphi$ חד־חד ערכית מהשאלה הקודמת, והטלת מנה היא תמיד על, לכן כדי לבדוק דיוק נותר לבדוק $\im \varphi = \ker \pi_{\im \varphi}$, אך זה ברור מהגדרה.
	נותר לבדוק פיצול, אבל זה מתקבל ישירות על־ידי $\psi$ מהשאלה הקודמת.
\end{proof}

\question{}
יהי $f : A \to B$ הומומורפיזם חוגים, ויהי $M$ מודול מעל $A$.

\subquestion{}
נראה שאם $M$ נוצר סופית אז $M_B$ נוצר סופית מעל $B$.
\begin{proof}
	$M$ נוצר סופית ולכן נסמן $S \subset M$ כך ש־$\langle S \rangle_A = M$ וכך ש־$S$ סופית.
	נזכור כי $M_B = B \otimes_A M$, מספיק להראות שכל טנזור פשוט נוצר סופית מעל $B$.
	יהי $b \otimes m \in M_B$, אז $m = \sum_{s \in S} a_s s$ עבור $a_s \in A$.
	נקבל,
	\[
		b \otimes m
		= b \otimes \sum a_s s
		= \sum_{s \in S} a_s (b \otimes s)
		= \sum_{s \in S} (f(a_s) b) (1 \otimes s)
	\]
	כאשר במעבר האחרון השתמשנו בהגדרת פעולת הכפל של $M_B$ כ־$B$־מודול.
	קיבלנו שניתן לייצג טנזור פשוט כצירוף לינארי של $1 \otimes s$ עבור $s \in S$ ולכן גם $M_B$ נוצר סופית.
\end{proof}

\subquestion{}
נראה שאם $M$ חופשי, אז גם $M_B$ חופשי מעל $B$.
\begin{proof}
	מהסעיף הקודם נותר להראות שאם $S$ בסיס אז גם $\{ 1 \otimes s \mid s \in S \}$ בסיס של $M_B$.
	מהגדרה שקולה לחופשיות נובע $M \cong \bigoplus_{i \in I} A$ עבור איזשהו $I$, ונזהה אותם.
	לכן,
	\[
		M_B
		= B \otimes_A M
		= B \otimes_A \bigoplus_{i \in I} A
		\cong \bigoplus_{i \in I} B \otimes_A A
		\cong \bigoplus_{i \in I} B
	\]
	וזהו בדיוק מודול חופשי מעל $B$.
\end{proof}

\question[7]{}
תהי אלגברה $A$ ממימד סופי מעל שדה $F$.

\subquestion{}
נראה שאם $A$ תחום שלמות אז $A$ שדה.
\begin{proof}
	יהי $0 \ne a \in A$, נראה כי הוא הפיך כפלית ולכן $A$ שדה.
	נניח כי $\dim_F A = n$ ונבחן את $\{ 1, a, a^2, \ldots, a^n \}$, זוהי קבוצה תלויה לינארית ולכן קיים צירוף לינארי כך שמתקיים,
	\[
		1
		= \sum_{i = 1}^n \alpha_i a^i
		= a \cdot \sum_{i = 1}^n \alpha_i a^{i - 1}
	\]
	ולכן נסמן $b = \alpha_1 1 + \cdots + \alpha_n a^{n - 1}$.
	נובע כי $1 = a \cdot b$, ו־$A$ תחום שלמות לכן $b \ne 0$, ונובע כי $a$ אכן הפיך.
\end{proof}

\subquestion{}
נראה שכל אידאל ראשוני של $A$ הוא מקסימלי.
\begin{proof}
	נרצה להראות שאם $I \triangleleft A$ אז גם $A / I$ אלגברה מעל $F$.
	ברור כי $A / I$ חוג חילופי, וברור כי $\alpha \cdot (a + I) = \alpha a + I$ מגדיר הומומורפיזם, לכן $A / I$ אלגברה מעל $F$.
	אנו יודעים כי קיימים $a_1, \ldots, a_n \in A$ פורשים של $A$, ולכן גם $a_1 + I, \ldots, a_n + I \in A / I$ בהכרח פורשים את $A / I$, כלומר $A / I$ ממימד סופי מעל $F$.
	נניח ש־$I$ ראשוני, לכן $A / I$ תחום שלמות, לכן הוא שדה מסעיף א', אז $I$ מקסימלי.
\end{proof}

\question{}
נראה שהחוג $\CC \otimes_{\RR} \CC$ אינו תחום שלמות.
\begin{proof}
	תחילה נבחין שזהו חוג כי $\CC$ היא אלגברה מעל $\RR$, לכן נוכל להגדיר את המכפלה הטנזורית של האלגברות כחוג על־ידי הפעולה,
	\[
		\left(\sum b_i \otimes c_i\right)\left(\sum b_j' \otimes c_j'\right)
		= \sum_{i, j} b_i b_j' \otimes c_i c_j'
	\]
	מהשאלה הקודמת נובע שמספיק להראות ש־$\CC \otimes_{\RR} \CC$ הוא לא שדה, כלומר שיש איבר לא הפיך.
	\[
		\varphi : \CC \times \CC \to \CC \otimes_{\RR} \CC,
		\qquad
		\varphi(z, w)
		= (z \otimes w) \cdot (1 \otimes i + i \otimes 1)
	\]
	במטלה הקודמת מצאנו ש־$1 \otimes i + i \otimes 1$ לא ניתן להצגה כטנזור פשוט, ונוכל להסיק שזהו לא איבר טריוויאלי, וכמוהו גם $1 \otimes i - i \otimes 1$, אבל,
	אבל,
	\[
		(1 \otimes i + i \otimes 1)(1 \otimes i - i \otimes 1)
		= (1 \otimes -1) - (i \otimes i) + (i \otimes i) - (-1 \otimes 1)
		= 0
	\]
	ולכן קיבלנו מחלקי אפס.
\end{proof}

\question[11]{}
נגדיר $A = C([0, 1], \CC)$, הפונקציות הרציפות $[0, 1] \to \CC$.
זוהי אלגברה מעל $\CC$, על־ידי כפל וחיבור פונקציות.

\subquestion{}
יהי $I \triangleleft A$ אידאל ממש.
נראה שקיים $x_0 \in [0, 1]$ כך שלכל $f \in I$ מתקיים $f(x_0) = 0$.
\begin{proof}
	נניח שהטענה לא מתקיימת ונראה ש־$c_1 \in I$, כלומר ש־$I = A$.
	יהי $x_0 \in [0, 1]$, ותהי $f \in I$ כך ש־$f(x_0) \ne 0$, לכן גם $\frac{f}{f(x_0)} \in I$, ובהתאם נניח ש־$f(x_0) = 1$.
	אילו $f$ בעלת שורש אז נוכל לקחת את $f + c_M$ עבור $M > \sup f$, ולכן נוכל להניח בלי הגבלת הכלליות ש־$f(x) \ne 0$ לכל $x \in [0, 1]$, ולכן $\frac{1}{f} \in A$ בהכרח, ונקבל $c_1 \in I$.
\end{proof}

\subquestion{}
נמצא את כל האידאלים המקסימליים של $A$.
\begin{proof}
	נגדיר את $M_{x_0} = \{ f \in A \mid f(x_0) = 0 \}$.
	אם $f, g \in M_{x_0}$ אז $(f + g)(x_0) = f(x_0) + g(x_0) = 0$ ולכן $f + g \in M_{x_0}$, וכן $c_0 \in M_{x_0}$ ולכן זהו אידאל.
	נניח ש־$M_{x_0} \subset I \triangleleft A$ אידאל, אז בהכרח מהסעיף הקודם קיימת נקודה $x_1$ כך ש־$f(x_1) = 0$ לכל $f \in I$.
	אם $x_0 = x_1$ אז נקבל $I = M_{x_0}$ בסתירה, ואם $x_1 \ne x_0$ אז בפרט נקבל שגם $f(x_1) = 0$ לכל $f \in M_{x_0}$ בסתירה להגדרתה, לכן $M_{x_0}$ אידאל מקסימלי.

	נסיק שהאוסף $\{ M_{x_0} \mid x_0 \in [0, 1] \}$ הוא בדיוק אוסף האידאלים המקסימליים של $A$.
\end{proof}

\subquestion{}
נראה של־$A$ יש אידאל ראשוני שאינו מקסימלי.
\begin{proof}
	נגדיר את הקבוצה,
	\[
		S = \{ f \in A \mid f(0) \ne 0, f(1) \ne 0 \}
	\]
	אז ברור כי $c_1 \in S$, וכן ש־$f g \in S$ לכל $f, g \in S$, לכן $S$ היא קבוצה כפלית.
	משאלה 13 במטלה 3 נובע כי קיים אידאל $P \triangleleft A$ ראשוני כך ש־$P \cap I = \varnothing$.
	נובע ש־$f(0) = f(1) = 0$ לכל $f \in P$, ולכן זהו לא אידאל מקסימלי.
\end{proof}

\end{document}
